% © 2025 Ing. David Jaroš — CC BY-NC-ND 4.0
%
% This work is licensed under a Creative Commons Attribution-NonCommercial-NoDerivatives
% 4.0 International License (CC BY-NC-ND 4.0).
%
% License History: Earlier drafts (up to v0.3) were released under CC BY 4.0.
% From v0.4 onward, all material is released under CC BY-NC-ND 4.0 to protect
% the integrity of the theoretical work during ongoing academic development.

\documentclass[12pt]{article}
\usepackage{amsmath,amssymb,amsthm}
\usepackage{mathtools}
\usepackage{geometry}
\usepackage{hyperref}
\geometry{margin=1in}

% Theorem environments
\newtheorem{definition}{Definition}[section]
\newtheorem{proposition}{Proposition}[section]
\newtheorem{theorem}{Theorem}[section]
\newtheorem{corollary}{Corollary}[section]
\newtheorem{remark}{Remark}[section]

% Custom commands
\newcommand{\R}{\mathbb{R}}
\newcommand{\calM}{\mathcal{M}}
\newcommand{\half}{\tfrac{1}{2}}

\title{%
  Phase 2: Multi-Scale Stability of the Effective Cosmological\\
  Parameter Manifold in Unified Biquaternion Theory\\[4pt]
  {\large Formal Derivation of the Rigidity Criterion}%
}
\author{Ing. David Jaroš}
\date{2025}

\begin{document}

\maketitle

\begin{abstract}
We formalise the numerical rigidity experiment that probes the structure of the
Layer-2 parameter manifold $\calM$ of Unified Biquaternion Theory (UBT).
Introducing an information-theoretic compatibility ratio
$p(\lambda)\in[0,1]$ at perturbation scale $\lambda\in\{0.8,1.0,1.2\}$ and its
associated surprise $I(\lambda)=-\log_2 p(\lambda)$, we define a quantitative
stability criterion: the manifold $\calM$ is \emph{rigid} if
$\Delta I \leq 2$\,bits and the hit-rate ratio $r \leq 3$ across all three scales.
We show that rigidity implies a constrained entropy manifold and a reduced volume
of viable cosmological parameter space, providing a sharp falsification target
for both UBT and $\Lambda$CDM.
\end{abstract}

\tableofcontents

%% =========================================================
\section{The Cosmological Parameter Manifold}
%% =========================================================

\subsection{Setup}

Let $\theta = (\theta_1,\dots,\theta_d)\in\calM\subset\R^d$ be the vector of
Layer-2 parameters (e.g.\ winding number $n_\psi$, RS code parameters $n,k$,
channel count $N_{\rm ch}$, and grid size $q$).
Fix a set of observable targets
$\mathbf{o} = (o_1,\dots,o_K)$ (e.g.\ fine-structure constant $\alpha$,
cosmological ratios) with experimental tolerances $\boldsymbol{\delta}$.

\begin{definition}[Hit condition]
\label{def:hit}
A configuration $\theta$ is a \emph{hit} at tolerance $\boldsymbol{\delta}$ if
\[
  \left|\frac{\hat{o}_j(\theta) - o_j}{o_j}\right| \leq \delta_j,
  \quad\forall j=1,\dots,K,
\]
where $\hat{o}_j(\theta)$ is the UBT prediction for observable $j$ at parameter
point $\theta$.
\end{definition}

\subsection{Perturbation Scaling}

For scale parameter $\lambda>0$ the perturbed sampling domain is
$\calM_\lambda = \{\theta : \theta_i\in[\mu_i - \lambda\sigma_i,\,\mu_i + \lambda\sigma_i]\}$,
where $\mu_i$ and $\sigma_i$ are the centre and half-width of the baseline range.
The three scales tested are
\begin{equation}
  \lambda \in \Lambda = \{0.8,\; 1.0,\; 1.2\},
  \label{eq:scales}
\end{equation}
corresponding to $-20\%$, baseline, and $+20\%$ range perturbations.

%% =========================================================
\section{Compatibility Ratio and Surprise}
%% =========================================================

\begin{definition}[Compatibility ratio]
\label{def:compat_ratio}
Given $N$ independent uniform samples $\{\theta^{(i)}\}_{i=1}^N$ drawn from
$\calM_\lambda$, the \emph{compatibility ratio} is the empirical hit rate:
\begin{equation}
  p(\lambda) = \frac{|\{i : \theta^{(i)}\;\text{is a hit}\}|}{N}.
  \label{eq:p_lambda}
\end{equation}
\end{definition}

\begin{definition}[Information surprise]
\label{def:surprise}
The \emph{information surprise} (rarity in bits) at scale $\lambda$ is:
\begin{equation}
  I(\lambda) = -\log_2 p(\lambda).
  \label{eq:I_lambda}
\end{equation}
If $p(\lambda)=0$ we set $I(\lambda)=+\infty$.  A large $I$ means the observational
targets are rare within the sampled parameter space.
\end{definition}

\begin{remark}
$I(\lambda)$ quantifies the number of binary questions one must ask to locate a
compatible configuration by random search.  It is an unbiased estimator of the
log-volume ratio $\log_2(\mathrm{Vol}(\calM_\lambda)/\mathrm{Vol}(\calM_\lambda^{\rm hit}))$.
\end{remark}

%% =========================================================
\section{Stability Criterion}
%% =========================================================

\subsection{Robustness Metrics}

For a completed three-scale experiment with results
$(p(0.8),\,p(1.0),\,p(1.2))$ define:

\begin{definition}[Max-delta]
\label{def:max_delta}
\begin{equation}
  \Delta I
  = \max_{\lambda\in\Lambda}\,|I(\lambda) - I(1.0)|
  \label{eq:Delta_I}
\end{equation}
\end{definition}

\begin{definition}[Hit-rate ratio]
\label{def:hit_ratio}
\begin{equation}
  r = \frac{\max_{\lambda\in\Lambda} p(\lambda)}{\min_{\lambda\in\Lambda} p(\lambda)}
  \label{eq:ratio_r}
\end{equation}
(with the convention $r=+\infty$ if $\min p(\lambda)=0$).
\end{definition}

\subsection{Formal Rigidity Condition}

\begin{definition}[Rigidity]
\label{def:rigidity}
The cosmological parameter manifold $\calM$ is \emph{rigid at scale $\Lambda$} if:
\begin{equation}
  \boxed{
    \Delta I \leq 2\;\text{bits}
    \quad\text{and}\quad
    r \leq 3.
  }
  \label{eq:rigidity}
\end{equation}
Otherwise $\calM$ is \emph{soft}.
\end{definition}

\begin{remark}
The thresholds $2$\,bits and $r=3$ correspond to at most a $4\times$ change in
hit-rate and at most a $4\times$ change in the rareness of a compatible solution
under $\pm 20\%$ range scaling.  Tighter choices would demand stronger stability;
looser choices would be physically uninformative.
\end{remark}

%% =========================================================
\section{Physical Interpretation}
%% =========================================================

\subsection{Rigidity and Constrained Entropy}

\begin{theorem}[Entropy constraint from rigidity]
\label{thm:entropy_constraint}
If $\calM$ is rigid (condition~\eqref{eq:rigidity}), then the differential
Shannon entropy $h(\theta)$ of a uniform distribution over $\calM_\lambda$ satisfies
\[
  |h(\calM_{1.2}) - h(\calM_{0.8})| \leq \log_2(3) \approx 1.58\;\text{bits}.
\]
\end{theorem}

\begin{proof}
$h(\calM_\lambda) = \log_2\mathrm{Vol}(\calM_\lambda)$.
Because the hit region scales as $\mathrm{Vol}(\calM_\lambda^{\rm hit})
= p(\lambda)\,\mathrm{Vol}(\calM_\lambda)$ and $r\leq 3$ bounds the ratio
$p(0.8)/p(1.2)\leq 3$,
\[
  h(\calM_{1.2}) - h(\calM_{0.8})
  = \log_2\frac{\mathrm{Vol}(\calM_{1.2})}{\mathrm{Vol}(\calM_{0.8})}
    + \log_2\frac{p(0.8)}{p(1.2)}
  \leq \log_2 3. \qedhere
\]
\end{proof}

\subsection{Rigidity and Reduced Viable Volume}

\begin{corollary}[Volume interpretation]
\label{cor:volume}
Rigidity implies that the \emph{fraction} of parameter space consistent with
observed cosmological constants changes by less than a factor of~$3$ when the
search range is varied by $\pm 20\%$.  Equivalently, the viable region of $\calM$
does not ``grow'' or ``shrink'' catastrophically under moderate range changes — a
signature of a \emph{constrained, not fine-tuned}, parameter space.
\end{corollary}

\subsection{Comparison with $\Lambda$CDM}

Standard $\Lambda$CDM has $\sim 6$ free parameters.  A randomly constructed
6-parameter space typically shows \emph{soft} behaviour: widening the prior
dramatically increases the compatible volume (large $r$, large $\Delta I$).
UBT's Layer-2 space has only $\sim 4$--6 parameters but they are algebraically
constrained (winding numbers, Reed--Solomon code parameters).
Rigidity of UBT's $\calM$ is therefore a \textbf{testable distinguishing
prediction}: it should be stable, whereas a generic $\Lambda$CDM prior should not.

\begin{proposition}[Falsification target]
\label{prop:falsif}
If the UBT parameter manifold is measured to be \emph{soft} ($r>3$ or $\Delta I>2$\,bits)
at any scale $\lambda$, this falsifies the claim that UBT's Layer-2 structure is
algebraically constrained rather than fine-tuned.
\end{proposition}

%% =========================================================
\section{Computational Protocol}
%% =========================================================

The experiment is implemented in
\texttt{forensic\_fingerprint/tools/layer2\_rigidity\_experiment.py}.
For each $\lambda\in\{0.8, 1.0, 1.2\}$:
\begin{enumerate}
  \item Sample $N$ configurations from $\calM_\lambda$ using a fixed random seed.
  \item Evaluate the hit condition (Definition~\ref{def:hit}).
  \item Record $p(\lambda)$ and $I(\lambda)$.
\end{enumerate}
Then compute $\Delta I$ and $r$, and apply Definition~\ref{def:rigidity} to emit
a \texttt{STABLE}/\texttt{UNSTABLE} verdict.  Reproducibility is guaranteed by
fixing the random seed; default \texttt{--seed~123}.

%% =========================================================
\section{Summary}
%% =========================================================

We have formalised the Layer-2 rigidity experiment as a well-defined mathematical
protocol.  The central objects are the compatibility ratio $p(\lambda)$
(Definition~\ref{def:compat_ratio}), the information surprise $I(\lambda)$
(Definition~\ref{def:surprise}), and the rigidity condition (Definition~\ref{def:rigidity}).
The physical interpretation (Section~4) shows that rigidity implies:
\begin{itemize}
  \item A constrained entropy manifold (Theorem~\ref{thm:entropy_constraint}).
  \item A reduced viable parameter volume (Corollary~\ref{cor:volume}).
  \item A testable distinction from $\Lambda$CDM (Proposition~\ref{prop:falsif}).
\end{itemize}
This section is ready for incorporation into a peer-reviewed manuscript.

\end{document}
